\documentclass[10pt, titlepage]{article}

\usepackage[margin=1in]{geometry}
\usepackage{fancyhdr}
\usepackage{amsfonts}
\usepackage{amsmath}
\usepackage{amssymb}
\usepackage{booktabs}
%\usepackage{tikz}
\usepackage{float}
\usepackage{bm}
\usepackage{pgfplots}


\title{\normalsize PLUS Session Worksheet \\[1em] \Large \textit{MATH 1554:} Linear Algebra}
\author{\textit{Covering Sections:} 2.4 \& 2.5}
\date{Fall 2023}

% ------- Set up header -------
\renewcommand{\sectionmark}[1]{\markright{#1}} % from https://tex.stackexchange.com/a/317393/227665
\pagestyle{fancy}
\fancyhead[L]{A. M., Fall 2023}
%\fancyhead[C]{\textit{\nouppercase{\rightmark}}}
\fancyhead[C]{\textit{Partition Matrices and  \(LU\) Factorization (2.2 \& 2.3)}}
\fancyhead[R]{\textbf{\thepage}}
% -------------------------------------------------------

% ------- Set up question environment -------
\newcounter{question}
\setcounter{question}{1}
\newenvironment{newquestion}
{
	\noindent\textbf{Question \thequestion}
}
{
\stepcounter{question}
\par
}
% -------------------------------------------------------

% ------- Set up Main Idea box environment -------
\newcounter{mainideas}
\setcounter{mainideas}{1}
\newcommand{\mainidea}{\textbf{Main Idea \#\themainideas \stepcounter{mainideas}}}
\setlength\fboxsep{1em}
%\newcommand{\midea}[2]{\begin{center}\fbox{\parbox{#1}{\centering \mainidea #2}} \end{center}} % Could be environment but I couldn't bother

% The way this works is that the savebox collects whatever was between the minipage and the end of it, then in the last line of the definition of the environment
% we regurgigate whatever was collected into an fbox which we center.
\newsavebox{\mybox}
\newenvironment{mideea}[1]
{%
\begin{lrbox}{\mybox}%
	\begin{minipage}{#1\textwidth}
		{\centering \mainidea \par}

}
{\end{minipage}\end{lrbox}\begin{center}\fbox{\usebox{\mybox}}\end{center}}
% -------------------------------------------------------


% ------- Set up math stuff -------
\renewcommand{\vec}[1]{\bm{#1}}
\DeclareMathOperator{\vspan}{Span}
\DeclareMathOperator{\row}{row}
\DeclareMathOperator{\col}{col}
\DeclareMathOperator{\colspace}{Col}
\DeclareMathOperator{\rowspace}{Row}
\DeclareMathOperator{\nullspace}{Nul}
\DeclareMathOperator{\rank}{rank}
\DeclareMathOperator{\proj}{proj}
% -------------------------------------------------------

% ------- Set up circle -------
\newcommand{\bubble}{\tikz[baseline=-0.625ex]\draw[radius=0.5em, fill=white] (0,0) circle;}
\newcommand{\dbubble}{\tikz[baseline=-0.625ex]\draw[radius=0.5em, fill=black] (0,0) circle;}
% -------------------------------------------------------

\makeatletter
\setlength{\@fptop}{0pt}
\makeatother

\begin{document}
\maketitle

\section{\centering T/F Questions}
\begin{newquestion}
	For each statement, mark the statement \emph{true} if it's \emph{always} true, or \emph{false} otherwise.
	\begin{table}[!h]
		\centering
		\begin{tabular}{c@{\hspace{2.5em}}cp{0.75\textwidth}}
			true & false & \\
			\specialrule{\lightrulewidth}{0.25em}{1.25em}
			\bubble & \bubble & If \(A\) is a \(n \times n\) matrix and \(A = LU\) is the \(LU\)-factorization of \(A\), then \(A\) is invertible if and only if \(U\) 
						      is.\\[0.75em]
			\bubble & \bubble & If \(A \in \mathbb{R}^{m \times n}\) and \(A = LU\) is its \(LU\)-factorization, then \(L \in \mathbb{R}^{n \times n}\) and
						      \(U \in \mathbb{R}^{m \times n}\).\\[0.75em]
			\bubble & \bubble & Every matrix has an \(LU\)-factorization.\\[0.75em]
			\bubble & \bubble & Every \emph{invertible} matrix has an \(LU\)-factorization.\\[0.75em]
			\bubble & \bubble & Given the matrices \(A, B, C\) and \(P = \begin{bmatrix}A & B \\ 0 & C\end{bmatrix}\) where \(P\) is \(p \times 5\) and 
						      \(A\) is \(n \times 3\), \(C\) is \(B\) is \((p-n) \times 2\).\\[0.75em]
			\bubble & \bubble & To get the \(LU\) factorization of any matrix \(A\), you must reduce \(A\) to RREF to get \(U\).\\[0.75em]
			\bottomrule
		\end{tabular}
	\end{table}
\end{newquestion}

\section{\centering Short Answer Questions}
\begin{newquestion}
	Answer the following questions.
	\begin{enumerate}
		\item Find the \(LU\) factorization of the matrix \(\begin{bmatrix}1 & 0 & -6 \\ 1 & 1 & -6 \\ 1 & 1 & -5\end{bmatrix}.\)

			\hfill \(L =  \fbox{\phantom{\(\begin{bmatrix}1 & 0 & -6 \\ 0 & 1 & 0 \\ 0 & 0 & 1\end{bmatrix}\)}}\)

			\hfill \(U = \fbox{\phantom{\(\begin{bmatrix}1 & 0 & 0 \\ -1 & 1 & 0 \\ 0 & -1 & 1\end{bmatrix}\)}}\)

		\item Use the \(LU\)-factorization obtained above to find a solution to the equation \(A\vec{x} = \begin{bmatrix}1 \\ 0 \\ 2\end{bmatrix}\).
		
		\item Given \(B, C, I_n \in \mathbb{R}^{n \times n}\) are all invertible and \(A = \begin{bmatrix}B & I_n \\ 0 & C\end{bmatrix}\), find \(X, Y, W, Z\)
			in terms of \(B, C,\) such that \(A^{-1} = \begin{bmatrix}X & Y \\ W & Z\end{bmatrix}\).
	\end{enumerate}
\end{newquestion}

\fancyhead[C]{\textit{Solutions}}

\section{\centering Solutions}
\subsection{Solution to T/F Questions}
	\begin{table}[!h]
		\centering
		\begin{tabular}{c@{\hspace{2.5em}}cp{0.75\textwidth}}
			true & false & \\
			\specialrule{\lightrulewidth}{0.25em}{1.25em}
			\dbubble & \bubble & If \(A\) is a \(n \times n\) matrix and \(A = LU\) is the \(LU\)-factorization of \(A\), then \(A\) is invertible if and only if \(U\) 
						       is.\\[0.75em]
			\bubble & \dbubble & If \(A \in \mathbb{R}^{m \times n}\) and \(A = LU\) is its \(LU\)-factorization, then \(L \in \mathbb{R}^{n \times n}\) and
						      \(U \in \mathbb{R}^{m \times n}\).\\[0.75em]
			\bubble & \dbubble & Every matrix has an \(LU\)-factorization.\\[0.75em]
			\bubble & \dbubble & Every \emph{invertible} matrix has an \(LU\)-factorization.\\[0.75em]
			\bubble & \dbubble & Given the matrices \(A, B, C\) and \(P = \begin{bmatrix}A & B \\ 0 & C\end{bmatrix}\) where \(P\) is \(p \times 5\) and 
						      \(A\) is \(n \times 3\), \(C\) is \(B\) is \((p-n) \times 2\).\\[0.75em]
			\bubble & \dbubble & To get the \(LU\) factorization of any matrix \(A\), you must reduce \(A\) to RREF to get \(U\).\\[0.75em]
			\bottomrule
		\end{tabular}
	\end{table}

\subsection{Solution to Short Answer Questions}
\begin{enumerate}
	\item \[L = \begin{bmatrix}1 & 0 & 0 \\ 1 & 1 & 0 \\ 1 & 1 & 1\end{bmatrix}, \qquad U = \begin{bmatrix}1 & 0 & -6 \\ 0 & 1 & 0 \\ 0 & 0 & 1\end{bmatrix}.\]
	\item First solve for \(L\vec{y} = \vec{b}\). You get that \(\vec{y} = \begin{bmatrix}1 \\ -1 \\ 2\end{bmatrix}\). Then, solve for \(U\vec{x} = \vec{y}\). You
		get that \(\vec{x} = \begin{bmatrix}13\\ -1 \\ 2\end{bmatrix}\).
	\item \[\begin{bmatrix}B & I_n \\ 0 & C\end{bmatrix}\begin{bmatrix}X & Y \\ W & Z\end{bmatrix} = \begin{bmatrix}BX + W & BY + Z \\ CW & CZ\end{bmatrix} = \begin{bmatrix}I_n & 0 \\ 0 & I_n\end{bmatrix}.\]
		We now have that \(CW = 0 \implies W = 0\), so \(BX + W = I_n \implies BX = I_n \implies X = B^{-1}.\) We then have that \(CZ = I_n \implies Z = 
		C^{-1}\). Finally, \(BY + Z = 0 \implies BY = -C^{-1} \implies Y = -B^{-1}C^{-1}.\)
\end{enumerate}


\end{document}